% !PRINTED 545
\weberpage{545}
Wir sprechen ihn so aus:

1. \gesperrt{Ist $\pi$ ein Primtheiler der natürlichen Primzahl $p$, so ist für jede ganze Zahl $\omega$ im Körper $\Omega$}
\[ \omega^{N_a(\pi)} \equiv \omega \pmod{\pi}. \]

Hieran knüpfen sich nun wichtige Folgerungen:

2. \gesperrt{Bezeichnet $f(t)$ eine ganze Function $m^{\text{ten}}$ Grades, deren Coëfficienten ganze Zahlen in $\Omega$ sind, und $\pi$ ein Primfunctional, so hat die Congruenz}
\weq{5}{f(t) \equiv 0 \pmod{\pi}}
\gesperrt{höchstens $m$ Wurzeln, d. h. es giebt höchstens $m$ incongruente ganze Zahlen in $\Omega$, die, für $t$ gesetzt, die Congruenz befriedigen.}

Bedeutet nämlich $\alpha$ irgend eine Zahl in $\mathfrak{o}$, so können wir
\weq{6}{f(t) = (t-\alpha)\,f_1(t) + f(\alpha)}
setzen, worin $f_1(t)$ eine ebensolche Function wie $f(t)$ ist, aber nur vom $(m-1)^{\text{ten}}$ Grade. Ist aber $f(\alpha) \equiv 0 \pmod{\pi}$, so muss jede Wurzel von (5) der Congruenz
\[ (t-\alpha)\,f_1(t) \equiv 0 \pmod{\pi} \]
genügen. Sie muss also entweder mit $\alpha$ congruent oder eine Wurzel der Congruenz $(m-1)^{\text{ten}}$ Grades
\[ f_1(t) \equiv 0 \pmod{\pi} \]
sein. Setzen wir unseren Satz als bewiesen voraus für Congruenzen $(m-1)^{\text{ten}}$ Grades, so gilt er demnach auch für Congruenzen $m^{\text{ten}}$ Grades; und da er für Congruenzen ersten Grades gilt, so ist er allgemein richtig.

Jede durch $\pi$ nicht theilbare Zahl in $\mathfrak{o}$ genügt, wie wir gesehen haben, der Congruenz
\weq{7}{\omega^{p^f-1} \equiv 1 \pmod{\pi}.}
Ist nun $a$ die kleinste natürliche Zahl, für die die Congruenz
\weq{8}{\omega^a \equiv 1 \pmod{\pi}}
befriedigt ist, so lässt sich durch das schon oft angewandte Schlussverfahren zeigen, dass jeder andere Exponent $l$, für den $\omega^l \equiv 1$ ist, ein Vielfaches von $a$ sein muss. Denn wäre $l$ nicht durch $a$ theilbar, so wäre auch, wenn $a'$ der Rest der Division von $a$ durch $l$ ist, $\omega^{a'} \equiv 1$, was nach der Voraussetzung über $a$ nicht möglich ist. % [sic] the print has ``von $a$ durch $l$''; the minimality argument requires the remainder of $l$ by $a$
Also ist $a$ ein Theiler von $p^f-1$, und wir nennen $\omega$ eine \gesperrt{zum Exponenten $a$ gehörige Zahl}.
